\DocumentMetadata{
  pdfversion=2.0,
  pdfstandard=ua-2,
  lang=en-US,
  tagging=on,
  tagging-setup={math/setup=mathml-SE,tabsorder=structure}
}
\documentclass[11pt,letterpaper]{article}
\usepackage[margin=0.68in,includefoot]{geometry}
\usepackage{newcomputermodern}
\usepackage{amsmath,amscd,mathtools}
\usepackage{tikz,adjustbox}
\providecommand{\square}{\mdlgwhtsquare}
\usepackage[hidelinks]{hyperref}
\hypersetup{
  pdftitle={Combinatorics PhD Exam, January 2018},
  pdfauthor={University of Florida Department of Mathematics},
  pdfsubject={Graduate mathematics examination},
  pdfdisplaydoctitle=true,
  bookmarksopen=true
}
\setcounter{secnumdepth}{0}
\setlength{\parindent}{0pt}
\setlength{\parskip}{0.28em}
\setlength{\emergencystretch}{3em}
\setlength{\leftmargini}{2.15em}
\setlength{\leftmarginii}{2.65em}
\setlength{\leftmarginiii}{3em}
\pagestyle{empty}
\raggedbottom
\allowdisplaybreaks
\NewTaggingSocketPlug{block/list/label}{examproblem}{%
  \tagstructbegin{tag=Lbl}\tagstructend%
  \tagstructbegin{tag=\LBody}%
  \tagstructbegin{tag=\ProblemHeadingTag,title-o={\ProblemHeadingText}}%
  \tagmcbegin{tag=\ProblemHeadingTag,actualtext={\ProblemHeadingText}}%
  #2%
  \tagmcend%
  \tagstructend%
}
\newenvironment{examproblems}[2]{%
  \def\ProblemHeadingTag{#2}%
  \AssignTaggingSocketPlug{block/list/label}{examproblem}%
  \begin{enumerate}%
  \setcounter{enumi}{#1}%
  \renewcommand{\labelenumi}{\textbf{\theenumi.}}%
  \setlength{\topsep}{0.35em}%
  \setlength{\partopsep}{0pt}%
  \setlength{\itemsep}{0.48em}%
  \setlength{\parsep}{0.18em}%
}{\end{enumerate}}
\newenvironment{examsubparts}{%
  \AssignTaggingSocketPlug{block/list/label}{default}%
  \begin{enumerate}%
  \setlength{\topsep}{0.18em}%
  \setlength{\partopsep}{0pt}%
  \setlength{\itemsep}{0.16em}%
  \setlength{\parsep}{0.1em}%
}{\end{enumerate}}
\newenvironment{examinstructions}{%
  \par\begingroup\itshape
}{%
  \par\endgroup\vspace{0.18em}
}
\makeatletter
\renewcommand\subsection{\@startsection{subsection}{2}{0pt}%
  {-1.25ex plus -.25ex minus -.1ex}%
  {0.55ex plus .1ex}%
  {\normalfont\large\bfseries}}
\renewcommand{\@maketitle}{%
  \newpage\null\vskip -1em%
  {\footnotesize\bfseries
    \href{https://www.ufl.edu/}{University of Florida}, \href{https://math.ufl.edu/}{Department of Mathematics}\par}%
  \vskip -0.5em%
  \tagstructbegin{tag=H1,title={Combinatorics PhD Exam, January 2018}}%
  \tagpdfparaOff%
  \tagmcbegin{tag=H1}%
    {\LARGE\bfseries\href{https://gma.math.ufl.edu/exams/combinatorics-phd/}{Combinatorics PhD Exam}, January 2018\par}%
  \tagmcend%
  \tagpdfparaOn%
  \tagstructend%
  \vskip 0.25em%
  \tagmcbegin{artifact}%
    \hrule height 1pt%
  \tagmcend%
  \vskip 1em%
}
\makeatother
\title{Combinatorics PhD Exam, January 2018}
\author{}
\date{}
\begin{document}
\maketitle
\thispagestyle{empty}
\begin{examproblems}{0}{H2}
\def\ProblemHeadingText{Problem 1.}
\item
Let \(D(n)\) be the number of derangements of length~\(n\). Let \(F(n)\) be the number of permutations of length~\(n\) that have exactly one fixed point. Determine (with proof) the number \(D(n)-F(n)\).
\def\ProblemHeadingText{Problem 2.}
\item
Prove: If~\(n\) and~\(m\) are the number of vertices and edges in a simple planar graph and \(n \ge 3\), then \(m \le 3n-6\).

\par
Prove: For a simple graph~\(G\) with \(n \ge 11\), at most one of~\(G\) or its complement \(\overline{G}\) is planar. (Recall that \(\overline{G}\) has the same vertex set as~\(G\) and, for all vertices \(u,v\), we have that \(\{u,v\}\) is an edge of \(\overline{G}\) if and only if \(\{u,v\}\) is not an edge of~\(G\).)
\def\ProblemHeadingText{Problem 3.}
\item
How many ways are there to seat~\(n\) married couples at a straight table (only one side of the table) so that no woman sits next to her husband?

\par
How many ways are there to seat~\(n\) married couples at a straight table so that men and women alternate and no woman sits next to her husband?
\def\ProblemHeadingText{Problem 4.}
\item
Let~\(B\) be a rectangular 3-dimensional box for which the width is equal to the length, but the depth is not equal to the length. How many ways are there to color the vertices of~\(B\) using colors red and blue, where two colorings are considered the same if one can be obtained from the other by an orientation preserving symmetry of~\(B\).
\def\ProblemHeadingText{Problem 5.}
\item
Consider a directed graph obtained by putting a direction on each edge of the complete graph. Such a digraph is called a tournament. Prove that a tournament has a Hamiltonian path. Is it true that every tournament has a Hamiltonian cycle? Explain?
\def\ProblemHeadingText{Problem 6.}
\item
Let \(n \ge 3\). Let \(h_n\) be the number of permutations of length~\(n\) in which the number of inversions is divisible by three. Find an explicit formula for \(h_n\).
\def\ProblemHeadingText{Problem 7.}
\item
Prove that the language \(\mathcal{L}=\{ww : w \in \{a,b\}^*\}\) is not regular.
\def\ProblemHeadingText{Problem 8.}
\item
Suppose~\(\mathcal{C}\) and~\(\mathcal{D}\) are permutation classes. Prove that \(\mathcal{C}\cup\mathcal{D}\) is a permutation class, and describe (with an algorithm, perhaps) how to compute its basis.
\end{examproblems}
\end{document}
